\documentclass[10pt,a4paper]{amsart}
\usepackage[margin=19mm]{geometry}
\usepackage{amsmath,amssymb,mathtools}
\usepackage[scr=rsfs,cal=euler]{mathalfa}
\usepackage{xfrac}
\usepackage[unicode,hidelinks,bookmarks=false]{hyperref}
\newcommand{\ie}{i.e.\ }
\newcommand{\cf}{cf.\ }
\newcommand{\resp}{respectively\ }
\newcommand{\Subsection}{\subsection}
\providecommand{\nobreakdash}{\nobreak\hbox{-}}
\setlength{\emergencystretch}{2em}
\multlinegap0pt
\usepackage{amssymb}
\usepackage{mathtools}
\newtheorem{prop}{Proposition}
\usepackage{cleveref}
\crefname{prop}{Proposition}{Propositions}
\newcommand{\Ho}{H}% symbol for Hopf algebra
\renewcommand{\P}{\mathcal{P}}
\newcommand{\R}{\mathcal{R}}
\newcommand{\cH}{\mathcal{H}}
\newcommand{\chabl}{\cH_{\lambda,\alpha,\beta}}
\newcommand{\ot}{\otimes}
\title{Deformation Theory and Hopf Actions on Koszul Algebras}
\author{A. V. Shepler, S. Witherspoon}
\date{Source: arXiv:2504.14397v3 (CC BY 4.0)}
\begin{document}
\maketitle

\noindent\emph{Source and licence.} Deformation Theory and Hopf Actions on Koszul Algebras. Version arXiv:2504.14397v3 (CC BY 4.0), distributed by arXiv under the Creative Commons Attribution 4.0 International licence, http://creativecommons.org/licenses/by/4.0/, as recorded in the arXiv licence metadata for this exact version and shown on the abstract page https://arxiv.org/abs/2504.14397v3. Full licence notice: \texttt{licenses/arxiv-2504.14397-CC-BY-4.0.txt}.\par\smallskip

\noindent\textit{Editorial scope.} Counted unit: the article's prop ChoiceParameterCodomain (source0.tex lines 5823--5842). The article's complete original proof of it is delivered with it (source0.tex lines 5843--5922, one \texttt{proof} environment of 80 lines). No mathematics is written, repaired or re-derived: the statement and the proof are the article's own lines, in the article's own notation, and no retained line is substituted or re-flowed.  Retained uncounted as the source-local context the proof needs: lines 1897--1899.  Nothing was omitted from any retained span, so the package carries no omission record.  No retained line contains a citation, so the wrapper carries no bibliography and no bibliography entry.  No retained line contains a figure environment, an image-inclusion command or drawing code, so no image is delivered and the package declares no asset dependency.  Editorial formatting only, in addition: an AMS \texttt{amsart} preamble that restates the article's own declarations and macros used by the retained lines, namely the environments \texttt{prop} on separate counters, together with the standard packages those lines need.  The retained environments print in this document as Proposition 1 in order of appearance, whereas the article prints the same items with the numbering of its own full document.  The complete native source of the article as distributed in the arXiv e-print for this exact version is bundled unchanged as \texttt{sources/176786/source0.tex}.
\begin{equation}\label{eqn:Hlab}
   \chabl \ =\ T_{\! _H}(\Ho\ot V \ot \Ho)/ (\P\cup \P')
\end{equation}

     \begin{prop}\label{ChoiceParameterCodomain}
      For any quotient algebra $\cH_{\lambda,\tilde\alpha, \tilde \beta}$
      defined by linear parameter functions
         %$\tilde\alpha: \R\rightarrow H\ot V\ot H$, $
     %\tilde\beta: \R\rightarrow H$, and $\lambda:\R'\rightarrow H$,
        $$\tilde\alpha: \R\rightarrow H\ot V\ot H, \qquad 
        \tilde\beta: \R\rightarrow H, \qquad\text{ and }\quad
        \lambda:\R'\rightarrow H, 
     $$
 there exist linear parameter functions 
     $\alpha: \R\rightarrow k1_H\ot V\ot H$ and 
     $\beta: \R\rightarrow H$
 %there exist linear parameter functions 
 %    $$\alpha: \R\rightarrow V\ot H, \qquad 
  %   \beta: \R\rightarrow H 
  %   \hspace{20ex}\hphantom{x}
  %   $$
with
$ \cH_{\lambda, \tilde\alpha, \tilde \beta} = \chabl$.
\end{prop}

\begin{proof}
    Consider the maps
    $$
    \begin{aligned}
     & \gamma:      H\ot V\ot H\longrightarrow 
     k1_H\ot V\ot \Ho \subset H\ot V \ot H,
     \quad\text{and}\quad 
     \\
     &\gamma': H\ot V\ot H\longrightarrow 
     \R'\cdot H \subset H\ot V \ot H
     \,
     \end{aligned}
     $$
     given by $\gamma(h\ot v\ot h') = (1_H\ot \tau(h\ot v))h'$
     and $\gamma'=\text{Id}-\gamma$.
We define parameter functions
     $$\alpha=\gamma\, \tilde\alpha:\ \R\longrightarrow k1_H\ot V\ot H
     \quad
     \text{ and }\quad
     \beta = \tilde \beta + \lambda \, \gamma' \, \tilde\alpha
     : \ \R\longrightarrow H
     $$
after extending $\lambda$ to 
     a right $H$-module homomorphism 
     $\lambda:H\ot V\ot H\rightarrow H$.

      Recall that $ \cH_{\lambda, \tilde\alpha, \tilde \beta}$ and $\chabl$
  are the quotients of 
  $T_H(H\ot V\ot H)$ by the ideals of filtered relations 
  $(\tilde{\P}\cup \P')$ and $(\P\cup \P')$, respectively, for 
         $$      
       \tilde{\P}=\{r-\tilde{\alpha}(r)-\tilde{\beta}(r):r\in \R\}
       \qquad\text{ and }\qquad 
       \P=\{r-\alpha(r)-\beta(r):r\in \R\}
       $$
       where 
       $\P'=\{r-\lambda(r):r\in \R'\}$
       (see~\cref{eqn:Hlab}). 
%
     To see that $ \cH_{\lambda, \tilde\alpha, \tilde \beta} =
       \chabl$, 
       we use the right $H$-module structure on $H\ot V \ot H$ and 
    argue that 
       ${\tilde\P}$ lies in the $k\text{-span}$ of $
     \{{\P}, 
     \P'\cdot H\}$
     whereas  ${\P}$ lies in $k\text{-span}
     \{\tilde{\P}, 
     \P'\cdot H\}$. 

     
     For $r$ in $\R$, abbreviate $x=\tilde{\alpha}(r)$ and
     $x'=\gamma'(x)$.  One may verify that
     $$\begin{aligned}
       r-\tilde{\alpha}(r)-\tilde{\beta}(r)
             & \ = \ r
       -\gamma(x)
       +\gamma(x) 
       -\tilde{\beta}(r)-\lambda(x')+\lambda(x') -x
       \\
       & \ = \ r-\alpha(r) +\gamma(x) 
       -{\beta}(r)+\lambda(x') -x
       \\
       &\ = \ r-\alpha(r) 
       -{\beta}(r)+\gamma(x)-x+\lambda(x')
       \\
       & \ = \ r-\alpha(r) 
       -{\beta}(r)-\gamma'(x)+\lambda(x')
       \ = \ (r-\alpha(r) 
       -{\beta}(r))-(x'-\lambda(x') )\, 
     \end{aligned}
     $$
       which lies in the $k\text{-span}
     \{{\P}, 
     \P'\cdot H\}$.
     Vice versa, this shows
     ${\P}$ lies in $k\text{-span}
     \{\tilde{\P}, 
     \P'\cdot H\}$. 
    \end{proof}

\end{document}
